% GENERATED FILE. Edit canonical lessons, then run npm run generate:books.
% canonical-source-hash: fnv1a64:b1d04ed347ceca54
% canonical-lessons: RU-C457-upakovyvat, RU-C457-zavorachivat, RU-C457-dostavka, RU-C457-dostavlyat, RU-C457-kuryer

\chapter{Packing and delivery}
\label{ch:ru-packing-and-delivery}
\hlchaptermodality{457}

\begin{chapteropening}
\textbf{By the end of this chapter:} \emph{I can send things: packing, wrapping, delivery, delivering and the courier.}

Complete the last lesson of chapter 457: I can send things: packing, wrapping, delivery, delivering and the courier.
\end{chapteropening}

\section[\texorpdfstring{\ru{упаковывать} (upakóvyvat)}{упаковывать (upakóvyvat)}]{\ru{упаковывать} (upakóvyvat) — to pack, to wrap}

\label{lesson:RU-C457-upakovyvat}



\begin{quote}
Before the new one: say the Russian for faulty, then the Russian for to save.
\end{quote}

\subsection*{You'll want to know: \ru{упаковывать}}
\textbf{\ru{упаковывать}} (\emph{upakóvyvat}) — \textquotedblleft{}to pack, to wrap\textquotedblright{}.

\begin{quote}
\ru{Продавец} \ru{упаковывает} \ru{подарок} \ru{в} \ru{красивую} \ru{бумагу}.
\end{quote}

\begin{quote}
\emph{Prodavéts upakóvyvayet podárok v krasívuyu bumágu.}
\end{quote}

\begin{quote}
The seller wraps the present in nice paper.
\end{quote}

Imperfective; perfective \emph{\ru{упаковать}} (\emph{\ru{упакую}, \ru{упакуешь}}).

\begin{cousinweb}
\emph{\ru{упаковывать}} is built from \emph{\ru{у}-} and \emph{\ru{паковать}} 'to pack', which is usually traced to German \emph{packen}. EtymDB does not list the verb, so the line is hedged.
\end{cousinweb}

\begin{grammarlens}[title={What you've built}]
One more word for telling, explaining and describing at B1.
\end{grammarlens}

\subsection*{Your turn}
\begin{itemize}
\raggedright
  \item \emph{Say it:} \emph{upakóvyvat}
  \item \emph{Say it:} \emph{upakóvyvat}, once more
  \item \emph{Say it:} a sentence of your own with \emph{upakóvyvat}
  \item \emph{Recall:} say the Russian for faulty, then the Russian for to save, then say \emph{upakóvyvat} again
\end{itemize}

\begin{tcolorbox}[breakable,colback=teal!4,colframe=teal!35!black,title={Before you move on}]
What does \ru{упаковывать} mean? (\textquotedblleft{}To pack, to wrap\textquotedblright{}.) Say it once more.
\end{tcolorbox}
\section[\texorpdfstring{\ru{заворачивать} (zavoráchivat)}{заворачивать (zavoráchivat)}]{\ru{заворачивать} (zavoráchivat) — to wrap (up); to turn (into)}

\label{lesson:RU-C457-zavorachivat}



\begin{quote}
Before the new one: say the Russian for to save, then the Russian for to pack.
\end{quote}

\subsection*{You'll want to know: \ru{заворачивать}}
\textbf{\ru{заворачивать}} (\emph{zavoráchivat}) — \textquotedblleft{}to wrap (up); to turn (into)\textquotedblright{}.

\begin{quote}
\ru{Продавщица} \ru{заворачивает} \ru{рыбу} \ru{в} \ru{бумагу}.
\end{quote}

\begin{quote}
\emph{Prodavshchítsa zavoráchivayet rýbu v bumágu.}
\end{quote}

\begin{quote}
The shop assistant wraps the fish in paper.
\end{quote}

Imperfective; perfective \emph{\ru{завернуть}}. Also to turn: \emph{\ru{завернуть} \ru{за} \ru{угол}} — to turn the corner.

\begin{cousinweb}
\emph{\ru{заворачивать}} is built from \emph{\ru{за}-} and \emph{-\ru{ворачивать}}, on the root of \emph{\ru{вертеть}} 'to turn': to turn something round in paper. EtymDB traces \emph{\ru{вертеть}} to Proto-Slavic \emph{vĭrtěti}.
\end{cousinweb}

\begin{grammarlens}[title={What you've built}]
One more word for telling, explaining and describing at B1.
\end{grammarlens}

\subsection*{Your turn}
\begin{itemize}
\raggedright
  \item \emph{Say it:} \emph{zavoráchivat}
  \item \emph{Say it:} \emph{zavoráchivat}, once more
  \item \emph{Say it:} a sentence of your own with \emph{zavoráchivat}
  \item \emph{Recall:} say the Russian for to save, then the Russian for to pack, then say \emph{zavoráchivat} again
\end{itemize}

\begin{tcolorbox}[breakable,colback=teal!4,colframe=teal!35!black,title={Before you move on}]
What does \ru{заворачивать} mean? (\textquotedblleft{}To wrap (up); to turn (into)\textquotedblright{}.) Say it once more.
\end{tcolorbox}
\section[\texorpdfstring{\ru{доставка} (dostávka)}{доставка (dostávka)}]{\ru{доставка} (dostávka) — delivery}

\label{lesson:RU-C457-dostavka}



\begin{quote}
Before the new one: say the Russian for to pack, then the Russian for to wrap.
\end{quote}

\subsection*{You'll want to know: \ru{доставка}}
\textbf{\ru{доставка}} (\emph{dostávka}) — \textquotedblleft{}delivery\textquotedblright{}.

\begin{quote}
\ru{Доставка} \ru{по} \ru{городу} \ru{бесплатная}.
\end{quote}

\begin{quote}
\emph{Dostávka po górodu besplátnaya.}
\end{quote}

\begin{quote}
Delivery within the town is free.
\end{quote}

Feminine. \emph{\ru{Доставка} \ru{на} \ru{дом}} — home delivery.

\begin{cousinweb}
\emph{\ru{доставка}} is built on \emph{\ru{доставить}} 'to deliver', from \emph{\ru{до}-} 'up to' and \emph{\ru{ставить}} 'to put, to place', which is inherited from Proto-Slavic \emph{staviti} (EtymDB).
\end{cousinweb}

\begin{grammarlens}[title={What you've built}]
One more word for telling, explaining and describing at B1.
\end{grammarlens}

\subsection*{Your turn}
\begin{itemize}
\raggedright
  \item \emph{Say it:} \emph{dostávka}
  \item \emph{Say it:} \emph{dostávka}, once more
  \item \emph{Say it:} a sentence of your own with \emph{dostávka}
  \item \emph{Recall:} say the Russian for to pack, then the Russian for to wrap, then say \emph{dostávka} again
\end{itemize}

\begin{tcolorbox}[breakable,colback=teal!4,colframe=teal!35!black,title={Before you move on}]
What does \ru{доставка} mean? (\textquotedblleft{}Delivery\textquotedblright{}.) Say it once more.
\end{tcolorbox}
\section[\texorpdfstring{\ru{доставлять} (dostavlyát)}{доставлять (dostavlyát)}]{\ru{доставлять} (dostavlyát) — to deliver}

\label{lesson:RU-C457-dostavlyat}



\begin{quote}
Before the new one: say the Russian for to wrap, then the Russian for delivery.
\end{quote}

\subsection*{You'll want to know: \ru{доставлять}}
\textbf{\ru{доставлять}} (\emph{dostavlyát}) — \textquotedblleft{}to deliver\textquotedblright{}.

\begin{quote}
\ru{Пиццу} \ru{доставляют} \ru{за} \ru{полчаса}.
\end{quote}

\begin{quote}
\emph{Pítstsu dostavlyáyut za polchasá.}
\end{quote}

\begin{quote}
They deliver the pizza in half an hour.
\end{quote}

Imperfective; the perfective is \emph{\ru{доставить}} (\emph{\ru{доставлю}, \ru{доставишь}}). \emph{\ru{Доставлять} \ru{удовольствие}} — to give pleasure.

\begin{cousinweb}
\emph{\ru{доставлять}} is built from \emph{\ru{до}-} 'up to' and \emph{-\ru{ставлять}}, from \emph{\ru{ставить}} 'to put, to place', which is inherited from Proto-Slavic \emph{staviti} (EtymDB).
\end{cousinweb}

\begin{grammarlens}[title={What you've built}]
One more word for telling, explaining and describing at B1.
\end{grammarlens}

\subsection*{Your turn}
\begin{itemize}
\raggedright
  \item \emph{Say it:} \emph{dostavlyát}
  \item \emph{Say it:} \emph{dostavlyát}, once more
  \item \emph{Say it:} a sentence of your own with \emph{dostavlyát}
  \item \emph{Recall:} say the Russian for to wrap, then the Russian for delivery, then say \emph{dostavlyát} again
\end{itemize}

\begin{tcolorbox}[breakable,colback=teal!4,colframe=teal!35!black,title={Before you move on}]
What does \ru{доставлять} mean? (\textquotedblleft{}To deliver\textquotedblright{}.) Say it once more.
\end{tcolorbox}
\section[\texorpdfstring{\ru{курьер} (kuryér)}{курьер (kuryér)}]{\ru{курьер} (kuryér) — a courier}

\label{lesson:RU-C457-kuryer}



\begin{quote}
Before the new one: say the Russian for delivery, then the Russian for to deliver.
\end{quote}

\subsection*{You'll want to know: \ru{курьер}}
\textbf{\ru{курьер}} (\emph{kuryér}) — \textquotedblleft{}a courier\textquotedblright{}.

\begin{quote}
\ru{Курьер} \ru{привёз} \ru{посылку} \ru{утром}.
\end{quote}

\begin{quote}
\emph{Kuryér privyóz posýlku útrom.}
\end{quote}

\begin{quote}
The courier brought the parcel in the morning.
\end{quote}

Masculine.

\begin{cousinweb}
\emph{\ru{курьер}} is usually traced to French \emph{courrier} 'a messenger', which EtymDB takes from Italian \emph{corriere}, a word usually linked with Latin \emph{currere} 'to run' (Lewis \& Short).
\end{cousinweb}

\begin{grammarlens}[title={What you've built}]
One more word for telling, explaining and describing at B1.
\end{grammarlens}

\subsection*{Your turn}
\begin{itemize}
\raggedright
  \item \emph{Say it:} \emph{kuryér}
  \item \emph{Say it:} \emph{kuryér}, once more
  \item \emph{Say it:} a sentence of your own with \emph{kuryér}
  \item \emph{Recall:} say the Russian for delivery, then the Russian for to deliver, then say \emph{kuryér} again
\end{itemize}

\begin{tcolorbox}[breakable,colback=teal!4,colframe=teal!35!black,title={Before you move on}]
What does \ru{курьер} mean? (\textquotedblleft{}A courier\textquotedblright{}.) Say it once more.
\end{tcolorbox}
